% Nota independiente. Compilar con latexmk -pdf calculos_para_rogelio.tex
\documentclass[11pt,a4paper]{article}
\usepackage[T1]{fontenc}
\usepackage[utf8]{inputenc}
\usepackage[spanish,es-nodecimaldot,es-noshorthands]{babel}
\usepackage{lmodern,microtype}
\usepackage[margin=24mm]{geometry}
\setlength{\parindent}{0pt}
\setlength{\parskip}{5pt plus 1pt}
\setlength{\emergencystretch}{2em}
\title{CSA para NAUSICA: puntos pendientes}
\author{J. J. G\'omez-Cadenas\\
 {\small Donostia International Physics Center (DIPC)}}
\date{\today}
\begin{document}
\maketitle
Los cálculos se han realizado con apoyo de Codex, versión Astra, high.

El cálculo aplica a tu circuito una corriente de iones distribuidos uniformemente durante 10 ms, con una capacitancia externa de 1.25 pF y una capacitancia de realimentación de 3 fF. LTspice calcula el voltaje y el espectro de ruido. Después se promedian los voltajes durante 300 $\mu$s en los extremos del frame y se calcula su diferencia, corregida por descarga. El mismo filtro se aplica al espectro de ruido. Se obtiene un ENC de unos 40 electrones RMS y un umbral de 203 iones para $S/N>5$.

Para completar el estudio, necesitamos:

Confirmar los modelos de transistores y el espectro de ruido que deben usarse. La ejecución con los archivos que me mandaste da una contribución flicker nula. Necesitaría el espectro numérico de ruido de salida entre 0.1 Hz y 1 MHz para esta configuración, a 20 $^\circ$C, junto con los modelos y la versión de LTspice que has utilizado. Con esos datos, puedo comprobar el ENC con Astra.

El artículo describe el cálculo. Los scripts, los modelos y los resultados están en el repositorio.
\end{document}
